\section{Introduction}
\label{sec:intro}

Retrieval-augmented generation (RAG) has become the dominant strategy for
grounding large language model outputs in external
evidence~\citep{lewis2020rag, gao2024retrievalaugmented}. The approach
underpins production systems from enterprise search to clinical
decision support, where factual reliability is not merely desirable but
safety-critical. The engineering logic behind most RAG development is
straightforward: retrieve more relevant passages, and the generator will
produce more faithful answers.

\paragraph{The problem.}
This intuition conflates two distinct quantities. Retrieval quality is a
property of \emph{individual passages}---each chunk is scored against the
query independently. Faithfulness is a property of the \emph{answer}, which
emerges from the generator reasoning over the full concatenation of retrieved
evidence. Between these two quantities lies a variable that standard RAG
evaluation ignores: the internal coherence of the retrieved context as a set.
When multiple passages are stitched into a prompt, their mutual consistency
determines whether the generator can synthesize a grounded response or is
forced to bridge evidential gaps with parametric hallucination.

\paragraph{Real-world consequences.}
The gap between retrieval quality and faithfulness is not merely theoretical.
In domain-specific deployments---biomedical question answering, legal
research, financial analysis---a general-purpose encoder may retrieve
passages that are individually topical yet collectively incoherent. Our
experiments show that under such conditions, RAG can produce \emph{more}
hallucinations than zero-shot prompting without any retrieved context,
turning a safety mechanism into a liability.

\paragraph{Why existing methods fall short.}
The RAG community has invested heavily in retrieval-side improvements:
re-ranking with cross-encoders~\citep{nogueira2019passage}, context
compression~\citep{xu2023recomp}, and self-reflective
retrieval~\citep{asai2024selfrag}. These methods optimize passage-level
relevance or brevity. None directly measure or preserve the coherence of
the assembled context. Post-retrieval refinement---splitting low-scoring
passages into sub-chunks and replacing them with higher-scoring
fragments---exemplifies this blind spot. It improves local matching scores
while silently destroying the connective tissue between evidence segments.

\paragraph{Our key idea.}
We hypothesize that faithfulness depends on a \emph{meso-level} property
of the context: whether the retrieved passages fit together coherently
enough for the generator to reason over them as a unified evidence set. We
test this hypothesis through a controlled evaluation covering 6,050+ queries
across two datasets, two generator families, and four retrieval
configurations. The central empirical finding is what we call the
\emph{refinement paradox}: aggressive post-retrieval refinement---the kind
commonly deployed in production RAG systems that maximize per-passage
relevance---raises per-chunk similarity while reducing answer faithfulness
by up to 15.9 percentage points and introducing hallucination where there
was none. The paradox is not that refinement is inherently harmful, but that
indiscriminate refinement optimizes the wrong objective.

Figure~\ref{fig:overview} summarizes the study design and the core finding.
\input{../figures/overview}

\paragraph{Contributions.}
This paper makes four contributions, each supported by controlled experiments
across 6,050+ queries:
\begin{enumerate}
  \item \textbf{We demonstrate that retrieval quality and answer
        faithfulness can move in opposite directions.}
        On SQuAD, naive post-retrieval refinement improves per-chunk
        similarity while reducing faithfulness from 0.7995 to 0.6407 and
        introducing 10.0\% hallucination ($-$15.9~pp). On PubMedQA, RAG
        with a general-domain encoder increases hallucination by
        2.5$\times$ relative to zero-shot prompting. This challenges the
        foundational assumption that better retrieval yields more faithful
        generation.

  \item \textbf{We show that selective, coherence-aware refinement
        eliminates the paradox.}
        HCPC-v2 recovers SQuAD faithfulness to 0.7874 with 0\%
        hallucination by intervening on only 16.7\% of queries. The method
        is robust across a 7$\times$ range of refinement rates
        (7\%--48\%) without degradation, demonstrating that the key design
        principle is restraint, not algorithmic sophistication.

  \item \textbf{We establish that chunk size dominates faithfulness
        variance by orders of magnitude over prompt strategy.}
        ANOVA reveals chunk size as the only significant factor on SQuAD
        ($F = 19.72$, $p < 0.001$, $\eta^2 = 0.100$, Cohen's $d = 1.41$),
        accounting for 90.1\% of the between-factor variance across the
        three tested configuration parameters, while prompt strategy
        explains $< 0.1$\%. This overturns the common engineering priority
        of optimizing prompts before retrieval configuration.

  \item \textbf{We introduce CCS, a retrieval-time diagnostic that
        separates safe from dangerous contexts without generation.}
        CCS consistently distinguishes coherent retrieval sets (0\%
        hallucination, CCS = 1.0 on SQuAD) from incoherent ones (20\%
        hallucination, CCS = 0.57 on PubMedQA), providing a computationally
        cheap guard-rail for deployment decisions.
\end{enumerate}

\FloatBarrier
